%%=============================================================================
%% Samenvatting
%%=============================================================================

% TODO: De "abstract" of samenvatting is een kernachtige (~ 1 blz. voor een
% thesis) synthese van het document.
%
% Een goede abstract biedt een kernachtig antwoord op volgende vragen:
%
% 1. Waarover gaat de bachelorproef?
% 2. Waarom heb je er over geschreven?
% 3. Hoe heb je het onderzoek uitgevoerd?
% 4. Wat waren de resultaten? Wat blijkt uit je onderzoek?
% 5. Wat betekenen je resultaten? Wat is de relevantie voor het werkveld?
%
% Daarom bestaat een abstract uit volgende componenten:
%
% - inleiding + kaderen thema
% - probleemstelling
% - (centrale) onderzoeksvraag
% - onderzoeksdoelstelling
% - methodologie
% - resultaten (beperk tot de belangrijkste, relevant voor de onderzoeksvraag)
% - conclusies, aanbevelingen, beperkingen
%
% LET OP! Een samenvatting is GEEN voorwoord!

%%---------- Nederlandse samenvatting -----------------------------------------
%
% TODO: Als je je bachelorproef in het Engels schrijft, moet je eerst een
% Nederlandse samenvatting invoegen. Haal daarvoor onderstaande code uit
% commentaar.
% Wie zijn bachelorproef in het Nederlands schrijft, kan dit negeren, de inhoud
% wordt niet in het document ingevoegd.

\IfLanguageName{english}{%
\selectlanguage{dutch}
\chapter*{Samenvatting}
\lipsum[1-4]
\selectlanguage{english}
}{}

%%---------- Samenvatting -----------------------------------------------------
% De samenvatting in de hoofdtaal van het document

\chapter*{\IfLanguageName{dutch}{Samenvatting}{Abstract}}
Het onderzoek in deze bachelorproef gaat over de uitdaging om een technisch complex, bestaande .NET-rekenmodel bij ILVO toegankelijker en beter reproduceerbaar te maken voor wetenschappelijk onderzoek en data-analyse in Python. De kern van het probleem is dat de huidige architectuur, hoewel functioneel sterk, voor onderzoekers vaak als een “black box” ervaren wordt: de berekeningen zijn aanwezig, maar de onderliggende logica, aannames, dataflow en validatiestappen zijn niet altijd zichtbaar, traceerbaar en goed gedocumenteerd. Daardoor ontstaan misinterpretaties, documentatie-drift en vertraging bij de ontwikkeling van nieuwe analyses.

De centrale onderzoeksvraag was hoe een hybride aanpak kan zorgen voor meer transparantie, reproduceerbaarheid en samenwerking zonder de bestaande C\#-architectuur volledig te moeten vervangen. Aan de hand van een analyse van de bestaande situatie, een studie van documentatie-drift en een proof-of-concept werd onderzocht of Python, in combinatie met Python.NET en Jupyter-notebooks, een passende oplossing kan vormen. De aanpak combineert de stabiliteit van de bestaande .NET-omgeving met de flexibiliteit van Python voor onderzoek, experimentatie en validatie.

Uit het onderzoek blijkt dat een volledige migratie naar Python niet noodzakelijk of wenselijk is. De bestaande technische basis blijft waardevol, maar kan beter ondersteunen door een duidelijke, gecontroleerde toegangslayer voor onderzoekers. De proof-of-concept toont aan dat relevante onderdelen van het model succesvol kunnen worden benaderd vanuit Python, zonder de onderliggende .NET-implementatie direct te herstructureren. Hierdoor kunnen analyses beter worden herhaald, gedeeld en gecontroleerd.

De conclusie is dat de grootste uitdaging niet alleen technisch, maar ook organisatorisch en communicatief ligt: een hybride, goed gedefinieerde architectuur, duidelijke domeinmodellen, geautomatiseerde validatie en docs-as-code zijn essentieel om de kwaliteit van zowel de software als de onderzoeksresultaten te verbeteren. Deze aanpak biedt ILVO een pragmatische en duurzame weg naar meer transparantie, minder documentation drift en beter reproduceerbaar onderzoek.


%%%%%%%%%%%%%%%%%%%%%%%%%%%%%%%%%%%%%%%%%%%%%%%%%%%%%%%%%%%%%%%%%%%%%%%%%%%%%%%%%%%%%%%%%%%%%%%%%%%%%%%%%%%%%%%%%%%%%%%%%%%%%%%%%%%%%%%%%%%%%%%%%%%%%%%%%%%%%%%%%%%%%%%%%%%%%%%%%%%%%%%%%%%%%%%%%%%%%%%%%%%%%%%%%%%%%%%%%%%%%%%%%%%%%%%%%%%%%%%%%%%%%%%%%%%%%%%%%%%%%%%%%%%%%%%%%%%%%%%%%%%%%%%%%%%%%%%%%%%%%%%%%%%%%%%%%%%%%%%%%%%%%%%%%%%%%%%%%%%%%%%%%%%%%%%%%%%%%%%%%%%%%%%%%%%%%%%%%%%%%%%%%%%%%%%%%%%%%%%%%%%%%%%%%%%%%%%%%%%%%%%%%%%%%%%%%%%%%%%%%%%%%%%%%%%%%%%%%%%%%%%%%%%%%%%%%%%%%%%%%%%%%%
%% Bij ILVO, het Instituut voor Landbouw-, Visserij- en Voedingsonderzoek,                                                                                                                                                                                                                                                                                                                                                                                                                          %%
%% draaien veel onderzoeksprojecten rond complexe rekenmodellen.                                                                                                                                                                                                                                                                                                                                                                                                                                    %%
%% Deze modellen beschrijven hoe je emissies moet berekenen,                                                                                                                                                                                                                                                                                                                                                                                                                                        %%
%% of hoe bepaalde bedrijfsparameters zich verhouden tot een bepaald resultaat.                                                                                                                                                                                                                                                                                                                                                                                                                     %%
%% Het probleem: onderzoekers schrijven deze modellen op,                                                                                                                                                                                                                                                                                                                                                                                                                                           %%
%% maar IT-developers implementeren ze vervolgens in C\# en .NET.                                                                                                                                                                                                                                                                                                                                                                                                                                   %%
%% Zodra dat gebeurd is, zijn onderzoekers min of meer afhankelijk van de IT-afdeling voor elke kleine aanpassing of test.                                                                                                                                                                                                                                                                                                                                                                          %%
%% De code is als een zwarte doos. Ze zien hun formule niet terug,                                                                                                                                                                                                                                                                                                                                                                                                                                  %%
%% ze begrijpen niet precies wat er ondertussen gebeurt,                                                                                                                                                                                                                                                                                                                                                                                                                                            %%
%% en elke keer dat ze iets willen proberen, moet de IT-afdeling er aan te pas komen.                                                                                                                                                                                                                                                                                                                                                                                                               %%
%%                                                                                                                                                                                                                                                                                                                                                                                                                                                                                                  %%
%% Deze bachelorproef bestudeert hoe die kloof kan worden gedicht. Het gaat uit van documentatie bij ILVO, analyseert waar de wrijving zit, en beschouwt welke bestaande werkwijzen uit soortgelijke domeinen (literate programming, Jupyter-notebooks, documentatietools) daar nuttig voor kunnen zijn. Met die inzichten is een proof-of-concept gebouwd als proof-point: kunnen onderzoekers meer autonomie krijgen en meer inzicht in hoe hun model werkt, zonder dat de IT-kant eronder lijdt? %%
%%                                                                                                                                                                                                                                                                                                                                                                                                                                                                                                  %%
%% Het blijkt inderdaad mogelijk. Door beter inzicht in de scheiding tussen wat onderzoekers moeten begrijpen en wat de implementatie doet, kan zo'n samenwerking veel vloeiender worden ingericht. Er is niet alles opnieuw nodig, maar wel betere documentatie, betere structuur, en ruimte voor onderzoekers om voorzichtig te experimenteren.                                                                                                                                                   %%
%%%%%%%%%%%%%%%%%%%%%%%%%%%%%%%%%%%%%%%%%%%%%%%%%%%%%%%%%%%%%%%%%%%%%%%%%%%%%%%%%%%%%%%%%%%%%%%%%%%%%%%%%%%%%%%%%%%%%%%%%%%%%%%%%%%%%%%%%%%%%%%%%%%%%%%%%%%%%%%%%%%%%%%%%%%%%%%%%%%%%%%%%%%%%%%%%%%%%%%%%%%%%%%%%%%%%%%%%%%%%%%%%%%%%%%%%%%%%%%%%%%%%%%%%%%%%%%%%%%%%%%%%%%%%%%%%%%%%%%%%%%%%%%%%%%%%%%%%%%%%%%%%%%%%%%%%%%%%%%%%%%%%%%%%%%%%%%%%%%%%%%%%%%%%%%%%%%%%%%%%%%%%%%%%%%%%%%%%%%%%%%%%%%%%%%%%%%%%%%%%%%%%%%%%%%%%%%%%%%%%%%%%%%%%%%%%%%%%%%%%%%%%%%%%%%%%%%%%%%%%%%%%%%%%%%%%%%%%%%%%%%%%%


